% GENERATED FILE. Edit canonical lessons, then run npm run generate:books.
% canonical-source-hash: fnv1a64:097eb18a397dea47
% canonical-lessons: KA-C432-hattikku, KA-C432-talebagu, KA-C432-takararu, KA-C432-adacane, KA-C432-sukta, KA-R432-review-426-432, KA-R432-review-final-432

\chapter{Pressure and Objections}
\label{ch:ka-pressure-and-objections}
\hlchaptermodality{432}

\begin{chapteropening}
\textbf{By the end of this chapter:} \emph{I can say to suppress, to bow down, an objection, a hindrance and suitable.}

Complete the last lesson of chapter 432: I can say to suppress, to bow down, an objection, a hindrance and suitable.
\end{chapteropening}

\section[\texorpdfstring{\kn{ಹತ್ತಿಕ್ಕು} (hattikku)}{ಹತ್ತಿಕ್ಕು (hattikku)}]{\kn{ಹತ್ತಿಕ್ಕು} (hattikku) — to suppress; to stifle}

\label{lesson:KA-C432-hattikku}



\begin{quote}
Before the new one: say the Kannada for a formal resolution, then the Kannada for to include.
\end{quote}

\subsection*{You'll want to know: \kn{ಹತ್ತಿಕ್ಕು}}
\textbf{\kn{ಹತ್ತಿಕ್ಕು}} — \emph{hattikku} — \textquotedblleft{}to suppress; to stifle\textquotedblright{}.

\begin{quote}
\kn{ಅವಳು} \kn{ತನ್ನ} \kn{ಕೋಪವನ್ನು} \kn{ಹತ್ತಿಕ್ಕಿದಳು}.
\end{quote}

\begin{quote}
\emph{avaḷu tanna kōpavannu hattikkidaḷu.}
\end{quote}

\begin{quote}
She suppressed her anger.
\end{quote}

Past (I …): \kn{ಹತ್ತಿಕ್ಕಿದೆ}, \emph{hattikkide}.

Of feelings and of protests.

\begin{cousinweb}
\textbf{\kn{ಹತ್ತಿಕ್ಕು}} (\emph{hattikku}) is a Kannada verb whose origin is not settled: DEDR does not list it.
\end{cousinweb}

\begin{grammarlens}[title={What you've built}]
One more word for saying what happens if.
\end{grammarlens}

\subsection*{Your turn}
\begin{itemize}
\raggedright
  \item \emph{Say it:} \emph{hattikku}
  \item \emph{Say it:} \emph{hattikku}, once more
  \item \emph{Say it:} \emph{avaḷu tanna kōpavannu hattikkidaḷu.} — She suppressed her anger.
  \item \emph{Recall:} say the Kannada for a formal resolution, then the Kannada for to include, then say \emph{hattikku} again
\end{itemize}

\begin{tcolorbox}[breakable,colback=teal!4,colframe=teal!35!black,title={Before you move on}]
What does \kn{ಹತ್ತಿಕ್ಕು} mean? (\textquotedblleft{}To suppress; to stifle\textquotedblright{}.) Say it once more.
\end{tcolorbox}
\section[\texorpdfstring{\kn{ತಲೆಬಾಗು} (talebāgu)}{ತಲೆಬಾಗು (talebāgu)}]{\kn{ತಲೆಬಾಗು} (talebāgu) — to bow down; to give in}

\label{lesson:KA-C432-talebagu}



\begin{quote}
Before the new one: say the Kannada for to include, then the Kannada for to suppress.
\end{quote}

\subsection*{You'll want to know: \kn{ತಲೆಬಾಗು}}
\textbf{\kn{ತಲೆಬಾಗು}} — \emph{talebāgu} — \textquotedblleft{}to bow down; to give in\textquotedblright{}.

\begin{quote}
\kn{ಒತ್ತಡಕ್ಕೆ} \kn{ತಲೆಬಾಗಿ} \kn{ಕಂಪನಿ} \kn{ಸಂಬಳ} \kn{ಹೆಚ್ಚಿಸಿತು}.
\end{quote}

\begin{quote}
\emph{ottaḍakke talebāgi kampani sambaḷa heccisitu.}
\end{quote}

\begin{quote}
Giving in to pressure, the company raised pay.
\end{quote}

Past (I …): \kn{ತಲೆಬಾಗಿದೆ}, \emph{talebāgide}.

\kn{ತಲೆ} (head) + \kn{ಬಾಗು} (to bend). What one gives in to takes -\kn{ಗೆ}.

\begin{cousinweb}
\textbf{\kn{ತಲೆಬಾಗು}} (\emph{talebāgu}) is native Dravidian, built from \emph{tale}, \textquotedblleft{}head\textquotedblright{} (DEDR 3103), and \emph{bāgu}, \textquotedblleft{}to bend, to bow\textquotedblright{} (DEDR 5335).
\end{cousinweb}

\begin{grammarlens}[title={What you've built}]
One more word for saying what happens if.
\end{grammarlens}

\subsection*{Your turn}
\begin{itemize}
\raggedright
  \item \emph{Say it:} \emph{talebāgu}
  \item \emph{Say it:} \emph{talebāgu}, once more
  \item \emph{Say it:} \emph{ottaḍakke talebāgi kampani sambaḷa heccisitu.} — Giving in to pressure, the company raised pay.
  \item \emph{Recall:} say the Kannada for to include, then the Kannada for to suppress, then say \emph{talebāgu} again
\end{itemize}

\begin{tcolorbox}[breakable,colback=teal!4,colframe=teal!35!black,title={Before you move on}]
What does \kn{ತಲೆಬಾಗು} mean? (\textquotedblleft{}To bow down; to give in\textquotedblright{}.) Say it once more.
\end{tcolorbox}
\section[\texorpdfstring{\kn{ತಕರಾರು} (takarāru)}{ತಕರಾರು (takarāru)}]{\kn{ತಕರಾರು} (takarāru) — an objection; a dispute}

\label{lesson:KA-C432-takararu}



\begin{quote}
Before the new one: say the Kannada for to suppress, then the Kannada for to bow down.
\end{quote}

\subsection*{You'll want to know: \kn{ತಕರಾರು}}
\textbf{\kn{ತಕರಾರು}} — \emph{takarāru} — \textquotedblleft{}an objection; a dispute\textquotedblright{}.

\begin{quote}
\kn{ಈ} \kn{ಹೊಸ} \kn{ಯೋಜನೆಗೆ} \kn{ಯಾರಿಗೂ} \kn{ತಕರಾರು} \kn{ಇಲ್ಲ}.
\end{quote}

\begin{quote}
\emph{ī hosa yōjanege yārigū takarāru illa.}
\end{quote}

\begin{quote}
No one has any objection to this new scheme.
\end{quote}

\kn{ತಕರಾರು} \kn{ಮಾಡು}, to raise objections.

\begin{cousinweb}
\textbf{\kn{ತಕರಾರು}} (\emph{takarāru}) is usually traced to Arabic \emph{takrār}, \textquotedblleft{}repetition, altercation\textquotedblright{}, through Hindustani.
\end{cousinweb}

\begin{grammarlens}[title={What you've built}]
One more word for saying what happens if.
\end{grammarlens}

\subsection*{Your turn}
\begin{itemize}
\raggedright
  \item \emph{Say it:} \emph{takarāru}
  \item \emph{Say it:} \emph{takarāru}, once more
  \item \emph{Say it:} \emph{ī hosa yōjanege yārigū takarāru illa.} — No one has any objection to this new scheme.
  \item \emph{Recall:} say the Kannada for to suppress, then the Kannada for to bow down, then say \emph{takarāru} again
\end{itemize}

\begin{tcolorbox}[breakable,colback=teal!4,colframe=teal!35!black,title={Before you move on}]
What does \kn{ತಕರಾರು} mean? (\textquotedblleft{}An objection; a dispute\textquotedblright{}.) Say it once more.
\end{tcolorbox}
\section[\texorpdfstring{\kn{ಅಡಚಣೆ} (aḍacaṇe)}{ಅಡಚಣೆ (aḍacaṇe)}]{\kn{ಅಡಚಣೆ} (aḍacaṇe) — a hindrance; a disruption}

\label{lesson:KA-C432-adacane}



\begin{quote}
Before the new one: say the Kannada for to bow down, then the Kannada for an objection.
\end{quote}

\subsection*{You'll want to know: \kn{ಅಡಚಣೆ}}
\textbf{\kn{ಅಡಚಣೆ}} — \emph{aḍacaṇe} — \textquotedblleft{}a hindrance; a disruption\textquotedblright{}.

\begin{quote}
\kn{ಮಳೆಯಿಂದ} \kn{ಪಂದ್ಯಕ್ಕೆ} \kn{ಅಡಚಣೆ} \kn{ಆಯಿತು}.
\end{quote}

\begin{quote}
\emph{maḷeyinda pandyakke aḍacaṇe āyitu.}
\end{quote}

\begin{quote}
The rain disrupted the match.
\end{quote}

\kn{ಅಡಚಣೆ} \kn{ಮಾಡು}, to disturb.

\begin{cousinweb}
\textbf{\kn{ಅಡಚಣೆ}} (\emph{aḍacaṇe}) is built on native \emph{aḍacu}, \textquotedblleft{}to press in, to cram\textquotedblright{} (DEDR 63).
\end{cousinweb}

\begin{grammarlens}[title={What you've built}]
One more word for saying what happens if.
\end{grammarlens}

\subsection*{Your turn}
\begin{itemize}
\raggedright
  \item \emph{Say it:} \emph{aḍacaṇe}
  \item \emph{Say it:} \emph{aḍacaṇe}, once more
  \item \emph{Say it:} \emph{maḷeyinda pandyakke aḍacaṇe āyitu.} — The rain disrupted the match.
  \item \emph{Recall:} say the Kannada for to bow down, then the Kannada for an objection, then say \emph{aḍacaṇe} again
\end{itemize}

\begin{tcolorbox}[breakable,colback=teal!4,colframe=teal!35!black,title={Before you move on}]
What does \kn{ಅಡಚಣೆ} mean? (\textquotedblleft{}A hindrance; a disruption\textquotedblright{}.) Say it once more.
\end{tcolorbox}
\section[\texorpdfstring{\kn{ಸೂಕ್ತ} (sūkta)}{ಸೂಕ್ತ (sūkta)}]{\kn{ಸೂಕ್ತ} (sūkta) — suitable; appropriate}

\label{lesson:KA-C432-sukta}



\begin{quote}
Before the new one: say the Kannada for an objection, then the Kannada for a hindrance.
\end{quote}

\subsection*{You'll want to know: \kn{ಸೂಕ್ತ}}
\textbf{\kn{ಸೂಕ್ತ}} — \emph{sūkta} — \textquotedblleft{}suitable; appropriate\textquotedblright{}.

\begin{quote}
\kn{ಈ} \kn{ಕೆಲಸಕ್ಕೆ} \kn{ಅವನೇ} \kn{ಸೂಕ್ತ} \kn{ವ್ಯಕ್ತಿ}.
\end{quote}

\begin{quote}
\emph{ī kelasakke avanē sūkta vyakti.}
\end{quote}

\begin{quote}
He is the right person for this job.
\end{quote}

\kn{ಸೂಕ್ತವಾದ}, fitting.

\begin{cousinweb}
\textbf{\kn{ಸೂಕ್ತ}} (\emph{sūkta}) is borrowed from Sanskrit \emph{sūkta}, \textquotedblleft{}well said, fitting\textquotedblright{} (CDIAL 13545).
\end{cousinweb}

\begin{grammarlens}[title={What you've built}]
One more word for saying what happens if.
\end{grammarlens}

\subsection*{Your turn}
\begin{itemize}
\raggedright
  \item \emph{Say it:} \emph{sūkta}
  \item \emph{Say it:} \emph{sūkta}, once more
  \item \emph{Say it:} \emph{ī kelasakke avanē sūkta vyakti.} — He is the right person for this job.
  \item \emph{Recall:} say the Kannada for an objection, then the Kannada for a hindrance, then say \emph{sūkta} again
\end{itemize}

\begin{tcolorbox}[breakable,colback=teal!4,colframe=teal!35!black,title={Before you move on}]
What does \kn{ಸೂಕ್ತ} mean? (\textquotedblleft{}Suitable; appropriate\textquotedblright{}.) Say it once more.
\end{tcolorbox}
\section[(dialogue)]{Rules, choices and objections}

\label{lesson:KA-R432-review-426-432}



\begin{quote}
Say the Kannada for to get rid of and to disregard.
\end{quote}

\subsection*{Your turn}
Say these aloud:

\begin{itemize}
\raggedright
  \item to get rid of, to disregard, careful, a snag, adverse
  \item a ratio, per cent, to classify, to convert, adequate
  \item if ever, or else, in any case, applicability, overall
  \item to hope for, to look to, to propose, to suggest, a solution
  \item a state of affairs, a circumstance, to arise, a stipulation, a requirement
  \item to conclude, to confer, a dialogue, a formal resolution, to include
  \item to suppress, to bow down, an objection, a hindrance, suitable
\end{itemize}

\begin{tcolorbox}[breakable,colback=teal!4,colframe=teal!35!black,title={Before you move on}]
To get rid of? (\textbf{\kn{ತೊಲಗಿಸು}}, \emph{tolagisu}.) To disregard? (\textbf{\kn{ಕಡೆಗಣಿಸು}}, \emph{kaḍegaṇisu}.) Careful? (\textbf{\kn{ಜಾಗರೂಕ}}, \emph{jāgarūka}.) A snag? (\textbf{\kn{ತೊಡಕು}}, \emph{toḍaku}.) Adverse? (\textbf{\kn{ಪ್ರತಿಕೂಲ}}, \emph{pratikūla}.) A ratio? (\textbf{\kn{ಅನುಪಾತ}}, \emph{anupāta}.) Per cent? (\textbf{\kn{ಶೇಕಡ}}, \emph{śēkaḍa}.) To classify? (\textbf{\kn{ವರ್ಗೀಕರಿಸು}}, \emph{vargīkarisu}.) To convert? (\textbf{\kn{ಪರಿವರ್ತಿಸು}}, \emph{parivartisu}.) Adequate? (\textbf{\kn{ಸಮರ್ಪಕ}}, \emph{samarpaka}.) If ever? (\textbf{\kn{ಒಂದೊಮ್ಮೆ}}, \emph{ondomme}.) Or else? (\textbf{\kn{ಇಲ್ಲವೇ}}, \emph{illavē}.) In any case? (\textbf{\kn{ಹೇಗಿದ್ದರೂ}}, \emph{hēgiddarū}.) Applicability? (\textbf{\kn{ಅನ್ವಯ}}, \emph{anvaya}.) Overall? (\textbf{\kn{ಒಟ್ಟಾರೆ}}, \emph{oṭṭāre}.) To hope for? (\textbf{\kn{ಆಶಿಸು}}, \emph{āśisu}.) To look to? (\textbf{\kn{ಅಪೇಕ್ಷಿಸು}}, \emph{apēkṣisu}.) To propose? (\textbf{\kn{ಪ್ರಸ್ತಾಪಿಸು}}, \emph{prastāpisu}.) To suggest? (\textbf{\kn{ಸೂಚಿಸು}}, \emph{sūcisu}.) A solution? (\textbf{\kn{ಪರಿಹಾರ}}, \emph{parihāra}.) A state of affairs? (\textbf{\kn{ಪರಿಸ್ಥಿತಿ}}, \emph{paristhiti}.) A circumstance? (\textbf{\kn{ಸಂದರ್ಭ}}, \emph{sandarbha}.) To arise? (\textbf{\kn{ಉದ್ಭವಿಸು}}, \emph{udbhavisu}.) A stipulation? (\textbf{\kn{ನಿಬಂಧನೆ}}, \emph{nibandhane}.) A requirement? (\textbf{\kn{ಅವಶ್ಯಕತೆ}}, \emph{avaśyakate}.) To conclude? (\textbf{\kn{ತೀರ್ಮಾನಿಸು}}, \emph{tīrmānisu}.) To confer? (\textbf{\kn{ಸಮಾಲೋಚಿಸು}}, \emph{samālōcisu}.) A dialogue? (\textbf{\kn{ಸಂವಾದ}}, \emph{saṃvāda}.) A formal resolution? (\textbf{\kn{ಠರಾವು}}, \emph{ṭharāvu}.) To include? (\textbf{\kn{ಒಳಗೊಳ್ಳು}}, \emph{oḷagoḷḷu}.) To suppress? (\textbf{\kn{ಹತ್ತಿಕ್ಕು}}, \emph{hattikku}.) To bow down? (\textbf{\kn{ತಲೆಬಾಗು}}, \emph{talebāgu}.) An objection? (\textbf{\kn{ತಕರಾರು}}, \emph{takarāru}.) A hindrance? (\textbf{\kn{ಅಡಚಣೆ}}, \emph{aḍacaṇe}.) Suitable? (\textbf{\kn{ಸೂಕ್ತ}}, \emph{sūkta}.)
\end{tcolorbox}
\section[(dialogue)]{Pressure and Objections — once more}

\label{lesson:KA-R432-review-final-432}



\begin{quote}
Say the Kannada for to suppress and to bow down.
\end{quote}

\subsection*{Your turn}
\begin{itemize}
\raggedright
  \item \emph{Say it:} to suppress, to bow down, an objection, a hindrance, suitable
\end{itemize}

\begin{tcolorbox}[breakable,colback=teal!4,colframe=teal!35!black,title={Before you move on}]
To suppress? (\textbf{\kn{ಹತ್ತಿಕ್ಕು}}, \emph{hattikku}.) To bow down? (\textbf{\kn{ತಲೆಬಾಗು}}, \emph{talebāgu}.) An objection? (\textbf{\kn{ತಕರಾರು}}, \emph{takarāru}.) A hindrance? (\textbf{\kn{ಅಡಚಣೆ}}, \emph{aḍacaṇe}.) Suitable? (\textbf{\kn{ಸೂಕ್ತ}}, \emph{sūkta}.)
\end{tcolorbox}
